\documentclass[11pt]{article}
\usepackage[T1]{fontenc}
\usepackage{amsmath,amssymb,amsthm,booktabs}
\usepackage[margin=1.06in]{geometry}
\usepackage[hidelinks]{hyperref}
\setlength{\emergencystretch}{2em}
\hypersetup{pdftitle={Eighth-root phases and strong unextendibility of mutually unbiased bases in dimension ten},pdfauthor={}}
\newtheorem{theorem}{Theorem}
\newtheorem{proposition}[theorem]{Proposition}
\newtheorem{lemma}[theorem]{Lemma}
\newtheorem{corollary}[theorem]{Corollary}
\newcommand{\C}{\mathbb C}
\newcommand{\Z}{\mathbb Z}
\newcommand{\Q}{\mathbb Q}
\newcommand{\BH}{\mathrm{BH}}
\newcommand{\ii}{\mathrm i}
\newcommand{\code}[1]{\texttt{\detokenize{#1}}}
\title{Eighth-root phases and strong unextendibility\\of mutually unbiased bases in dimension ten}
\author{}
\date{}
\begin{document}
\maketitle
\vspace{-2.2em}
\begin{abstract}
Fix an orthonormal coordinate basis of \(\C^{10}\). Suppose that two further
bases are mutually unbiased to it and to one another, and that within every
vector of the two bases all coordinate ratios are eighth roots of unity.
We prove that no complex unit vector is unbiased to all three bases.
The finite part of the proof classifies octonary unbiased mates for each
of the 60 monomial equivalence classes of \(\BH(10,8)\). Only two
representatives admit mates, with 36 and one unordered dephased mates,
respectively. A common change of coordinate phases then makes both bases
quaternary up to row phases. An exact Gaussian-integer Laurent identity
excludes an additional vector with arbitrary complex phases. The
classification coverage, mate enumeration, identity, and passage to
complex measurement probabilities are verified in Lean~4. Consequently,
in a complete set of eleven bases, after fixing the coordinate vectors of
one basis, at least nine of the remaining ten would have a vector with a
coordinate ratio outside the eighth roots of unity.
\end{abstract}

\section{The obstruction}

Two orthonormal bases of \(\C^d\) are \emph{mutually unbiased} if
\(\lvert\langle u,v\rangle\rvert^2=1/d\) for all vectors \(u,v\) from
different bases. A family of such bases is \emph{strongly unextendible}
if no unit vector has squared overlap \(1/d\) with every basis vector in
the family. Whether \(\C^{10}\) admits a complete set of eleven mutually
unbiased bases remains open \cite{yang}. We give a condition on two bases
in one common coordinate system that already rules out an additional
unbiased vector.

Write \(\mu_q=\{z\in\C:z^q=1\}\). Fix an ordered orthonormal basis
\(E=(e_0,\ldots,e_9)\), including a choice of phase for each \(e_j\).
Write \(u_j=\langle e_j,u\rangle\), with the inner product linear in its
second argument. A unit vector unbiased to \(E\) satisfies
\(\lvert u_j\rvert=1/\sqrt{10}\) at every coordinate.
We say that its ray has \emph{eighth-root coordinate ratios} if
\[
                    (u_j/u_0)^8=1 \qquad(0\le j<10).
\]
This condition is invariant under an arbitrary phase on \(u\).
The coordinate phases of \(E\) are shared by all rays under consideration.

\begin{theorem}\label{th:main}
Let \(E,H,K\) be pairwise mutually unbiased bases of \(\C^{10}\).
If each of the twenty rays of \(H\) and \(K\) has eighth-root coordinate
ratios in \(E\), the triple \(E,H,K\) is strongly unextendible.
\end{theorem}

The restriction on phases concerns the given bases. The theorem quantifies
over the whole unit sphere for the candidate vector.

\begin{corollary}\label{co:eleven}
Suppose that eleven pairwise mutually unbiased bases exist in \(\C^{10}\).
Choose any member as \(E\) and fix its coordinate phases. At least nine
of the ten other bases contain a vector \(u\) and a coordinate \(j\)
for which \((u_j/u_0)^8\ne1\).
\end{corollary}
\begin{proof}
If two of those ten bases had the eighth-root property for every ray,
Theorem~\ref{th:main} would apply. Any vector from a fourth member of
the family would contradict the theorem.
\end{proof}

A Butson Hadamard matrix \(A\in\BH(n,q)\) has entries in \(\mu_q\)
and satisfies \(AA^*=nI_n\). Best and Kharaghani gave an explicit unbiased
pair in \(\BH(10,4)\) and proved an upper bound of two mutually unbiased
quaternary matrices at orders twice an odd integer \cite{best}.
Their bound concerns an entire additional matrix with quaternary entries.
Kharaghani, Tayfeh-Rezaie and Zaitsev subsequently counted the
quaternary mates of each order-ten quaternary representative
\cite[Table~4]{regular}. Single quaternary matrices of order ten were
classified in \cite{quaternary}, and Lampio, \"Osterg{\aa}rd and
Sz\"oll\H{o}si classified the 60 single octonary matrices up to
monomial equivalence \cite{orderly,data}. We use the published
representatives and independently certify their coverage. The mate
classification below determines which of these matrices have an
unbiased octonary mate and which phases occur in that mate. The
final obstruction concerns one further vector over \(\C\).

Yang, Liu, Li and Deng proved strong unextendibility for triples of
product bases of \(\C^2\otimes\C^5\), and for triples containing a pair
equivalent to the standard and Fourier bases \cite[Theorems~A and~B]{yang}.
Here the hypothesis is a restriction on coordinate ratios in a common
basis. The proof proceeds by simultaneous reduction to four phases
and the quaternary certificate of the companion software \cite{software}.
It uses neither the Fourier-pair exclusion nor the product-basis exclusion.

\section{Phases and exact overlaps}

Set \(\zeta=\exp(\pi\ii/4)\), so that \(\zeta^2=\ii\).
Multiplying every coordinate of a normalized flat row by \(\sqrt{10}\)
makes each entry unimodular. We use these \emph{raw} rows for all
finite computations. For raw rows \(h,r\), put
\[
                  S(h,r)=\sum_{j=0}^{9}\overline{h_j}r_j .
\]
The corresponding Born probability is \(\lvert S(h,r)\rvert^2/100\).
Consequently distinct rows in a basis have \(S=0\), and rows in two
mutually unbiased bases have \(\lvert S\rvert^2=10\).
Every ray in the hypothesis of Theorem~\ref{th:main} has a unique raw
representative with first entry \(1\) and all entries in \(\mu_8\).

If \(h_j=\zeta^{p_j}\) and \(r_j=\zeta^{q_j}\), their overlap is a
sum of ten eighth roots with exponents \(q_j-p_j\bmod 8\). Group
the even and odd exponents to write
\begin{equation}\label{eq:split}
          S=A+\zeta B,\qquad A=a+b\ii,\quad B=u+v\ii ,
          \qquad a,b,u,v\in\Z .
\end{equation}
Each integer coefficient lies in \([-10,10]\).
The four directions \(1,\ii,\zeta,\ii\zeta\) are linearly independent
over \(\Q\), since \(x^4+1\) is the minimal polynomial of \(\zeta\).

\begin{lemma}\label{le:overlap}
The overlap is zero exactly when \(a=b=u=v=0\). It has squared modulus
\(10\) exactly when one of \(A,B\) is zero and the absolute values of
the real and imaginary parts of the other are \(1,3\) in either order.
\end{lemma}
\begin{proof}
The zero criterion follows from linear independence. For the norm, put
\(x=a^2+b^2\) and \(y=u^2+v^2\). Direct expansion gives
\[
   |A+\zeta B|^2
     =x+y+\sqrt{2}\bigl(a(u-v)+b(u+v)\bigr).
\]
If this equals \(10\), irrationality of \(\sqrt2\) gives \(x+y=10\)
and \(au+bv=av-bu\). The two-square identity gives
\(xy=(au+bv)^2+(av-bu)^2=2(au+bv)^2\).
Up to interchange, \(x=0,1,\ldots,5\). The possible products \(xy\)
are \(0,9,16,21,24,25\), of which only \(0\) is twice an integer square.
A Gaussian integer of norm ten has squared components \(1\) and \(9\).
The converse follows directly from \eqref{eq:split}.
\end{proof}

For an efficient exact test define
\[
 L(a,b,u,v)=a+32b+32^2u+32^3v .
\]
On the coefficient box just described this map is injective. Reducing
an equality modulo 32 identifies the first signed coefficient, because
the difference of the two coefficients has absolute value at most 20.
Subtraction and division by 32 recover the remaining coefficients
successively. The contribution of an exponent \(0,1,\ldots,7\) to \(L\)
is, in that order,
\[
                1,\ 1024,\ 32,\ 32768,\ -1,\ -1024,\ -32,\ -32768 .
\]
Thus raw orthogonality has integer score zero. The acceptable scores
for unbiasedness form the sixteen-element set
\begin{equation}\label{eq:targets}
   T=U\cup 1024U,\qquad
   U=\{\pm97,\pm95,\pm35,\pm29\}.
\end{equation}
Both equivalences are verified against the ordinary complex inner product
in the formalization.

\section{Classifying the octonary mates}

A monomial equivalence of a single octonary matrix permits independent
row and column permutations and multiplication of individual rows and
columns by elements of \(\mu_8\). In a pair, a common column operation
acts on the ambient coordinates. Each basis retains its own row
permutations and row phases. We fix the first matrix \(H\) at a listed
representative. A \emph{mate} of \(H\) is an unordered set of ten
distinct octonary rows, dephased to have first entry 1, that are mutually
orthogonal and are each unbiased to every row of \(H\) in the raw
convention. The count keeps \(H\) fixed and retains mates related by
its stabilizer whenever they give different sets of dephased rows.

\begin{proposition}[Finite classification]\label{pr:scan}
In the order of the published Aalto file \code{BH-10-8.txt}, numbering
its lines from one, only representatives 9 and 28 have an octonary
mate. They have 36 and one distinct mates, respectively. These two
representatives, as supplied, and every dephased row in each of their
mates have all exponents even.
\end{proposition}

\begin{proof}
We describe the exhaustive searches and the checks on their output.

\emph{Single-matrix coverage.}
Dephase a matrix so that its first row and first column are all ones.
Use exponents in \(\{0,\ldots,7\}\) and number coordinates from zero.
There is a representative whose rows and columns are simultaneously
lexicographically nondecreasing. To see this, take a minimum of
\(\sum_{r,j}8^{9-r}8^{9-j}p_{rj}\) over permutations fixing the first
row and column of a dephased exponent matrix \(p\). The base-eight
value of a row orders the rows lexicographically. Swapping an inverted
pair decreases the displayed integer. The same argument applies to
columns by symmetry. The first row and column consist of zeros, so
they do not introduce inversions.

Build such a representative in increasing row order, beginning with
the zero exponent row. The second row is nondecreasing. Its occurrences
of \(t\) and \(t+4\) must balance, by Lemma~\ref{le:overlap}. If
\(n_0,\ldots,n_3\) denote the four pair counts, then
\(\sum n_t=5\) and \(n_0\ge1\) because the first entry is zero.
This gives \(\binom73=35\) possible second rows. At each subsequent
step require orthogonality to the preceding rows. Whenever adjacent
columns have agreed throughout the prefix, the appended entries in
these columns must appear in nondecreasing order. These tests retain
every prefix of a doubly ordered solution. The remaining pool must
also contain enough distinct rows to finish the matrix. The resulting
search yields \(15\,564\) ten-row leaves. For every leaf a witness
specifies a column permutation and row and column phases taking its row
multiset to a listed representative. The verified list of leaves is
equal to the list carrying these witnesses.

The code checks the Gram matrix of each representative. It also
separates the 60 representatives using an integer invariant of the
exponent matrix. For clarity, write
\(d_{abkj}=[p_{bj}-p_{aj}-p_{bk}+p_{ak}]_8\), where brackets denote
least nonnegative residues. The invariant is
\[
 \sum_{a,b,c,d,k}
 \left[\sum_j(d_{abkj}+d_{ackj}+d_{adkj})^2\right]_{19}.
\]
Row and column phases cancel in the four-entry differences. Permutations
reindex the sums. Direct evaluation gives 60 distinct values. This
certifies inequivalence under the stated monomial action.

\emph{Rows for a fixed representative.}
Construct all \(8^9\) possible dephased rows implicitly. Split a row
into its first five and last five entries. The first half has
\(8^4\) choices and the second has \(8^5\). For the first two
anchor rows, tabulate the pair of partial scores of each second half.
For each first half with partial scores \((s_1,s_2)\), and each
\((t_1,t_2)\in T^2\), retrieve all second halves with scores
\((t_1-s_1,t_2-s_2)\). Additivity of \(L\) ensures that this finds
every solution of the first two equations. Filter against the remaining
anchor rows using \eqref{eq:targets}. Shared initial anchor rows can
reuse the same partial tables. Table~\ref{tab:rays} gives the resulting
numbers of distinct dephased rays. The sum over these sixty fixed
representatives is \(11\,200\).

\begin{table}[tb]
\centering
\caption{Octonary rays unbiased to the standard basis and to each
Aalto representative. Within each row the entries follow consecutive
one-based catalogue indices. Bold entries mark anchors admitting a
complete mate.}\label{tab:rays}
\small
\begin{tabular}{@{}lrrrrrrrrrr@{}}
\toprule
Indices & \multicolumn{10}{c}{Number of distinct dephased rays}\\
\midrule
1--10  &560&216&400&560&224&336&320&304&\textbf{912}&400\\
11--20 &304&232&144&160&192&128&32&60&128&192\\
21--30 &232&144&232&144&192&128&172&\textbf{112}&360&80\\
31--40 &360&128&24&48&60&224&224&400&400&320\\
41--50 &144&232&192&192&192&24&224&24&32&24\\
51--60 &20&22&22&0&8&20&20&20&0&0\\
\bottomrule
\end{tabular}
\end{table}

\emph{Bases from rows.}
In the graph on the surviving rays join two vertices when their
integer overlap score is zero. A mate is precisely a ten-clique.
Recursively choosing one vertex, intersecting the remaining pool with
its neighbours and ordering the choices lists every ten-clique.
An independently stored list of rays and a list of unordered mate
sets supply positive witnesses. Lean checks their validity, their
cardinalities as sets, and that every valid ray or mate occurs in
the corresponding list. The mate counts are 36 at index 9, one at
index 28, and zero at the other 58 indices. Checking the two anchor
tables and every row of the surviving cliques gives the asserted
parity of the exponents. Individual surviving rays that do not belong
to a ten-clique are not subject to this conclusion.
\end{proof}

\begin{corollary}[Simultaneous reduction]\label{co:descent}
Under the hypotheses of Theorem~\ref{th:main}, one common monomial
coordinate change and separate phases on the twenty rows give raw
matrices \(H',K'\in\BH(10,4)\). In particular, for the stacked raw
matrix \(X\) of twenty rows,
\[
                  \frac{X_{aj}X_{bk}}{X_{ak}X_{bj}}\in\mu_4
           \qquad\text{for all rows \(a,b\) and coordinates \(j,k\)} .
\]
The two rows may belong to different bases.
\end{corollary}
\begin{proof}
Represent both bases by dephased eighth-root rows and move the first
basis to the catalogue by monomial equivalence. Apply the same column
operation to the second basis and dephase its rows again. The Gram
conditions are preserved, so its rows form a mate in
Proposition~\ref{pr:scan}. Both transformed matrices have only even
exponents. Row and column phases cancel in the displayed ratios, and a
column permutation merely relabels \(j,k\).
\end{proof}

\section{The obstruction for an arbitrary complex vector}

It remains to exclude a flat vector whose entries can range over the full
circle. We describe the characteristic-zero certificate for this step
and its transport to all the pairs found above. In the quaternary
catalogue of \cite{quaternary}, exhaustive row and clique enumeration
again gives 36 mates for one representative and one for another,
in agreement with \cite[Table~4]{regular}. The quaternary part of
\cite{software} additionally supplies explicit equivalence witnesses
taking every such pair to the pair in \cite[Table~1]{best}, by a common
coordinate transformation, row permutations, row phases and a possible
exchange of the two bases. Thus one certificate for those
twenty reference rows suffices. For completeness we spell out its
equations and the check that makes it an obstruction.
The quaternary proofs and certificate data used here are available at
\href{https://github.com/sansssssss1/mub10-lean/tree/4321a0e7c888149fbeadd5ab8356555a68baa64e}{revision \code{4321a0e}} of \cite{software}.

Write \(R_{rj}=\ii^{p_{rj}}\) for the reference rows
\(0\le r<20\). The normalized input vectors are \(R_r/\sqrt{10}\).
Any unit vector unbiased to \(E\) has the form \(v=z/\sqrt{10}\) with
\(|z_j|=1\). In the Laurent ring
\(\Z[\ii][z_0^{\pm1},\ldots,z_9^{\pm1}]\) form
\[
 f_r(z)=\sum_{j\ne k}\ii^{\,p_{rk}-p_{rj}}\,z_jz_k^{-1}.
\]
On the torus \(z_k^{-1}=\overline{z_k}\), so expansion of the square
gives
\begin{equation}\label{eq:laurent}
       f_r(z)=\left|\sum_j\overline{R_{rj}}z_j\right|^2-10
              =100\left(\left|\left\langle
                  R_r/\sqrt{10},v\right\rangle\right|^2-\tfrac1{10}\right).
\end{equation}
The diagonal terms give the subtracted ten. A common unbiased vector
would therefore make all twenty polynomials vanish.

\begin{lemma}\label{le:cert}
For the reference rows there are Gaussian-integer Laurent polynomials
\(g_0,\ldots,g_{19}\) and a positive integer \(C\) such that
\(\sum_{r=0}^{19} g_r f_r=C\).
\end{lemma}
\begin{proof}
The witness in \cite{software} consists of \(24\,140\) triples
\((r,m,c)\), with \(m\in\Z^{10}\) and \(c\in\Z[\ii]\). For each row
index put \(g_r=\sum_{(r,m,c)}c z^m\).
To check the identity, initialize a dictionary indexed by vectors
in \(\Z^{10}\) with value zero. For each triple and each ordered
pair \(j\ne k\), add \(c\,\ii^{p_{rk}-p_{rj}}\) at key
\(m+e_j-e_k\), using \(e_j\) here for the \(j\)-th integer unit
vector. The \(2\,172\,600\) contributions occupy \(24\,691\) distinct
keys before cancellation. Collecting them leaves a single nonzero
coefficient at the zero exponent. It is the positive 313-digit integer \(C\)
specified by \code{constantValue} in \code{Certificate.lean}.
The finite identity \code{exact_certificate} verifies the collection
with exact signed integers and checks the bounds required by the
packed exponent representation. The evaluation lemma in
\code{StrongTriple.lean} then identifies that constant with the
linear combination of the complex row equations.
\end{proof}

One can reconstruct the witness from the reference matrices. The
reproduction script selects 610 independent orbit rows using reduction
modulo 13, solves the resulting system over the rationals, expands the
orbit multipliers, and clears denominators and a common factor. The
proof of Lemma~\ref{le:cert} uses the re-expanded integer identity
directly. As a numerical scale for this certificate, the sum of
\(|\Re c|+|\Im c|\) over its multipliers is strictly less than
\(44\,278 C\). On the unit torus this implies that some row has
\(|f_r|\ge 1/44\,278\), giving a probability discrepancy of at least
\(1/4\,427\,800\) for the reference pair by \eqref{eq:laurent}.
These integer and real inequalities are also checked in the quaternary
development \cite{software}.

\begin{proof}[Proof of Theorem~\ref{th:main}]
Apply Corollary~\ref{co:descent} and the quaternary equivalence witnesses
to obtain the reference pair. A common monomial coordinate change
preserves flatness, and applying the change to the state and to each
measurement vector preserves the absolute inner products. Row phases
and basis permutations preserve the requirement that every probability
equal \(1/10\). If a common unbiased vector existed, its transformed
coordinates would therefore give \(f_0=\cdots=f_{19}=0\) in
\eqref{eq:laurent}. Substitution into Lemma~\ref{le:cert} would force
the positive integer \(C\) to vanish.
\end{proof}

\section{Verification and scope}

The implementation uses Lean~4.19.0 and mathlib~4.19.0.
The equivalence of the integer scores and the complex overlap tests
is proved in \code{OctSixtyPhysical.lean}.
The file \code{OctSingleClassification.lean} proves coverage
for arbitrary \(\BH(10,8)\) inputs. It combines a proof of lexicographic
normalization, completeness of the row recursion, and a checked monomial
witness for each leaf. \code{OctDistinct.lean} distinguishes all sixty
representatives. The exact ray and mate sets are checked in
\code{OctSixtyExact.lean}. Candidate tables need satisfy these checks
before they are used in the physical theorem.

The file \code{OctSixtyConclusions.lean} contains both principal statements.
\begin{center}
\small
\begin{tabular}{@{}ll@{}}
Theorem~\ref{th:main} & \code{eight_relative_no_vector}\\
Corollary~\ref{co:eleven} & \code{complete_at_least_nine_violating_eighth_power}
\end{tabular}
\end{center}
The latter counts bases containing a ratio whose eighth power
differs from \(1\). The equivalence between the eight encoded phases and
all complex roots of \(z^8=1\) is proved in \code{OctAlphabet.lean}.
From the repository root, run the following commands to build the proofs
and independently recompute the finite certificates \cite{software}.
\begin{verbatim}
lake exe cache get
lake build
node reproduce/octonary.mjs --forward
\end{verbatim}
The last command regenerates the doubly ordered leaves and compares the
ray, mate, and leaf witnesses with the stored data. An independent exact
reconstruction of the Laurent identity is available in
\code{reproduce/recover.py}, with its Python dependency pinned in
\code{reproduce/requirements.txt}. The published single-matrix
dataset is available under CC BY 4.0 \cite{data}.

The dependency audit for the final theorems lists three standard
principles, \code{propext}, \code{Classical.choice}, and \code{Quot.sound}.
It also lists \code{Lean.ofReduceBool}, which records reliance on the native evaluator for finite Boolean
checks, including the search results and the collected Laurent identity.
Thus reproducing the verification uses the pinned Lean compiler and its
native evaluation path. Single-matrix catalogue coverage is itself a
theorem of the development.

Theorem~\ref{th:main} gives an obstruction for triples with a common
eighth-root coordinate description. A putative complete family could
satisfy Corollary~\ref{co:eleven} by using more general phases in nine
or all ten remaining bases. The unrestricted existence question for
dimension ten therefore remains open.

\paragraph{Computational assistance.}
OpenAI Codex assisted with the proof development, independent
recomputations and manuscript drafting. The formal checks reported here
were performed by Lean. Scientific Agent Skills \cite{skills} provided
procedural guidance for the writing and verification workflow.

\begin{thebibliography}{9}
\raggedright
\bibitem{best}
D.~Best and H.~Kharaghani,
Unbiased complex Hadamard matrices and bases,
\emph{Cryptogr. Commun.} \textbf{2} (2010), 199--209.
\href{https://doi.org/10.1007/s12095-010-0029-8}{doi:10.1007/s12095-010-0029-8}.

\bibitem{quaternary}
P.~H.~J.~Lampio, F.~Sz\"oll\H{o}si and P.~R.~J.~\"Osterg{\aa}rd,
The quaternary complex Hadamard matrices of orders 10, 12, and 14,
\emph{Discrete Math.} \textbf{313} (2013), 189--206.
\href{https://doi.org/10.1016/j.disc.2012.10.001}{doi:10.1016/j.disc.2012.10.001}.

\bibitem{orderly}
P.~H.~J.~Lampio, P.~R.~J.~\"Osterg{\aa}rd and F.~Sz\"oll\H{o}si,
Orderly generation of Butson Hadamard matrices,
\emph{Math. Comp.} \textbf{89} (2020), 313--331.
\href{https://doi.org/10.1090/mcom/3453}{doi:10.1090/mcom/3453}.

\bibitem{data}
P.~Lampio, P.~\"Osterg{\aa}rd and F.~Sz\"oll\H{o}si,
Dataset for Orderly Generation of Butson Hadamard Matrices, Zenodo (2019).
\href{https://doi.org/10.5281/zenodo.2585766}{doi:10.5281/zenodo.2585766}.

\bibitem{regular}
H.~Kharaghani, B.~Tayfeh-Rezaie and V.~Zaitsev,
On regular quaternary Hadamard matrices, \emph{J. Combin. Des.}
\textbf{34} (2026), 215--227.
\href{https://doi.org/10.1002/jcd.70006}{doi:10.1002/jcd.70006}.

\bibitem{yang}
J.~Yang, Y.~Liu, Y.~Li and X.-H.~Deng,
No complete set of mutually unbiased bases in dimension ten contains the
Fourier pair or three product bases (2026).
\href{https://arxiv.org/abs/2611.21188}{arXiv:2611.21188}.

\bibitem{software}
\emph{mub10-lean: Octonary and quaternary obstructions for mutually
unbiased bases in dimension ten}, Lean~4 software and certificates (2026).
\url{https://github.com/sansssssss1/mub10-lean}.

\bibitem{skills}
T.~Kassis, V.~Agarwal, Y.~He, D.~Patel and A.~M.~Brueckner,
Scientific Agent Skills: A Library of Procedural Knowledge for Research
Agents (2026).
\href{https://arxiv.org/abs/2609.00065}{arXiv:2609.00065}.
\end{thebibliography}
\end{document}
